% ============================================================
% ch28.tex —— 第 28 章 平行线惹的祸：非欧几何与宇宙的形状
% ============================================================
\chapter{平行线惹的祸：非欧几何与宇宙的形状}

欧几里得《几何原本》（约公元前 300 年）从五条公设出发，建起了两千年的几何大厦。前四条简洁得像诗：任意两点定一直线；线段可任意延长；给定中心与半径可作圆；所有直角彼此相等。第五条却啰嗦得扎眼：

\begin{quote}\itshape
平面内一条直线与另两条直线相交，若某一侧的两内角之和小于两直角，则这两条直线在该侧充分延长后必相交。
\end{quote}

这就是\textbf{第五公设}。它等价于初中课本那条"过直线外一点，\textbf{有且只有}一条直线与已知直线平行"（普莱费尔公理，1795 年改写的等价版）。两千年来人人觉得它是"疑似定理"：长得像能从前四条推出来的东西，只是还没找到证法。托勒密试过（证错）、普罗克鲁斯试过（偷偷用了待证结论）、中世纪阿拉伯学者集体攻关、十八世纪萨凯里差点成功（他推出了双曲几何的一堆定理，却认定它们"与直觉矛盾"而放弃——离诺贝尔级发现一步之遥，转身骂自己荒谬）、勒让德在教科书里反复返工——\textbf{两千年，全败}。

\section{惊天反转：它推不出来，因为它可以换}

十九世纪初，三个人先后做了同一件胆大包天的事。匈牙利的\textbf{鲍耶}（J\'anos Bolyai）、俄国的\textbf{罗巴切夫斯基}（1829 年发表）、以及更早想通却秘而不宣的\textbf{高斯}（私人信件里说"我怕愚人的嚷嚷"）。

操作是：\textbf{把第五公设反过来}——假设"过直线外一点可以作\textbf{不止一条}平行线"（具体说：有无穷多条），保留前四条公设，硬着头皮往下推，看逻辑会不会崩。

推了一条又一条：三角形内角和\textbf{小于} $180^\circ$；圆周长与测地半径之比\textbf{大于} $2\pi$；不存在相似而不全等的三角形；三角形的面积有上限……几百条定理次第成立，\textbf{没有一条撞出矛盾}。鲍耶在给父亲的信里写下那句著名的宣告：

\begin{quote}\centering\itshape
"我从无到有，创造了一个崭新的整个世界。"\quad（Aus dem Nichts habe ich eine seltsame neue Welt geschaffen.）
\end{quote}

逻辑不崩意味着什么？意味着\textbf{第五公设独立于前四条}——它不是定理，是\textbf{选择}。欧氏几何只是众多自洽几何中的一种。更妙的是：你在第 26、27 章已经提前认识了这个新世界——内角和小于 $180^\circ$、$C > 2\pi\rho$，这不正是\textbf{负曲率}世界（马鞍面/薯片）的测量结果吗！罗巴切夫斯基的"新世界"不是逻辑怪物，是\textbf{大薯片的几何}（严格说：完备的双曲平面无法完整放进三维空间，但它每片的局部都是薯片）。

黎曼（1854 年著名的就职演讲《论几何学基础中的假设》，哥廷根，听众里有年迈的高斯）补上了另一头：假设\textbf{没有平行线}（任何两条"直线"必相交）——那就是球面几何，内角和大于 $180^\circ$，$C < 2\pi\rho$，\textbf{正曲率}。黎曼更进一步：允许曲率\textbf{逐点变化}（不再是常数），并把舞台从曲面推广到任意维——广义相对论的数学（第 27 章蚂蚁、本章宇宙）在他演讲那天就已备好，只等六十年后爱因斯坦出生。

三种几何同桌而坐：

\begin{center}
\begin{tabular}{lccc}
\toprule
 & 双曲几何 & 欧氏几何 & 球面几何 \\
\midrule
过线外一点的平行线 & 无穷多条 & 恰好一条 & 一条都没有 \\
内在曲率 $K$ & $<0$（薯片） & $=0$（平面） & $>0$（球） \\
三角形内角和 & $<180^\circ$ & $=180^\circ$ & $>180^\circ$ \\
圆周长 $/$ 测地半径 & $>2\pi$ & $=2\pi$ & $<2\pi$ \\
三角形面积 & 有上限（角亏封顶） & 无上限 & 有上限（半球封顶） \\
相似三角形 & 不存在（角定则形定） & 存在 & 不存在 \\
\bottomrule
\end{tabular}
\end{center}

\begin{faming}
"几何"从此是\textbf{复数}。三种几何在逻辑上\textbf{同样自洽}——数学无法裁决谁"真"，裁决权移交\textbf{实验}：我们的宇宙三内角和到底是多少？这个问题问数学家是问错了门，要问物理学家。
\end{faming}

\section{实验的判决：宇宙的形状}

宇宙是哪种几何？这是观测问题，而人类已经量到了。两路证据：

\textbf{路线一：宇宙三角学。}找一个大三角形量内角和。人类手头最大的三角形：观测到的\textbf{宇宙微波背景辐射}（CMB，大爆炸的余晖，约 $138$ 亿光年外的"最后散射面"）。普朗克卫星（2018 年数据）通过 CMB 温度涨落的相关图案反推宇宙的曲率：$\Omega_K = 0.001 \pm 0.002$——\textbf{可观测宇宙的曲率与零相差不到千分之二}（误差内与 $0$ 相容）。欧氏几何统治我们能看到的每一寸空间；但"极其接近零"不等于"是零"——若宇宙整体带一丝正弯，它就是一个半径至少几百亿光年的超巨型球面，第 26 章"操场量不出地球弯"的剧本在宇宙尺度重播：\textbf{我们可能永远在"操场上"}。

\textbf{路线二：物质弯时空。}爱因斯坦 1915 年的广义相对论走得更根本：\textbf{有物质的地方，时空就地弯曲}——曲率与物质密度由场方程直接挂钩。地球绕太阳，不是被"引力绳"拉着，而是沿弯曲时空里的\textbf{测地线}（四维的"绷紧线"，第 26 章）自由滑行；光线经过太阳附近也会"拐弯"（1919 年日食观测证实，爱因斯坦一夜成名）；GPS 卫星的时钟每天比地面快约 $38$ 微秒，不吃相对论修正，导航一天偏十公里——\textbf{弯曲时空是每天校准你手机定位的工程现实}。第 25 章的曲率、第 26 章的测地线、第 27 章的内在弯曲，在物理里全部上桌：\textbf{物质告诉时空如何弯曲，时空告诉物质如何运动}（惠勒的名句）。黎曼 1854 年的抽象演讲，六十年后成了物理学的官方语言。

\section{方法论大结局}

把两千年的公设悬案复盘成三步——你会发现它就是本书的整个剧本：

\begin{enumerate}
\item \textbf{质疑定义}：第五公设是"显然的真理"还是"被发明的设计"？——序章的发明许可证，用在了数学最神圣的文件上；
\item \textbf{发明替代品}：把公设换掉，规则重新自洽即可——第 7 章"体检合格就是群"的翻版：体检合格（自洽）就是几何；
\item \textbf{实验裁决}：数学管自洽，物理管真假——第 6 章 RSA（数论管安全）、第 23 章 $\mathrm{li}(x)$（分析管素数）都在说同一句话：\textbf{数学的地图不止一张，选哪张走，是地图之外的事}。
\end{enumerate}

\begin{lishi}
三位非欧几何先驱的命运各不相同。罗巴切夫斯基：喀山大学校长，1846 年被解职（政治原因），晚年失明，靠口述完成最后著作，生前几乎无人理解——今天喀山大学立有他的雕像，月球背面有以他命名的环形山。鲍耶·亚诺什：父亲鲍耶·法尔卡什（高斯的老友）写信警告他"别碰平行线，我比你还懂这深渊的可怕"；儿子偏要碰，1825 年左右做出成果，附在父亲 1832 年的书的附录里寄给高斯——高斯回信说"称赞它等于称赞我自己，这工作的全部内容你的儿子……与我几十年的思考惊人一致"。这句话\textbf{摧毁了鲍耶}：他以为遭了剽窃，从此再没发表数学。高斯本人：什么都没发。三人三种结局，一个定理——数学史的显微镜下，"优先权"三个字何其昂贵。
\end{lishi}

\begin{dongshou}
\begin{itemize}
  \yisi 在球面上给"没有平行线"补齐论证：第 26 章习题证过"任何两条大圆必相交于一对对径点"。用同一条思路证明球面上"过直线外一点作不出任何平行线"（其实更强：连一条都不存在）。
  \yier 双曲世界"三角形面积有上限"的直觉核算：角亏（$180^\circ -$ 内角和）正比于面积（吉拉德定理的镜像版），而内角和最低能压到逼近 $0^\circ$（三个角都可以几乎为零）——所以面积封顶在"角亏 $180^\circ$"对应的最大值。请用"角趋近零、面积趋近上限"的语言描述双曲世界里一个"巨大三角形"长什么样（三条超长边，三个尖角，面积却不再增长）。
  \yier 用地球经线的图像向家人解释"宇宙接近平直但未必是平的"：两条经线在赤道附近"近乎平行"，到两极汇合——\textbf{近看平行，不等于永远平行}；反过来，你眼前的桌面永远水平，也说服不了你"宇宙是平的"。把这段话讲给别人听（费曼技巧：讲得清才算懂）。
\end{itemize}
\end{dongshou}

\begin{shaonao}
\textbf{双曲世界不存在相似形}（角角角定全等）的证明骨架，请补全细节：

(1) 双曲几何中"角亏正比于面积"（吉拉德的镜像，系数 $|K|$）：$180^\circ - (\text{内角和}) = |K|\times\text{面积}$。

(2) 假设存在两个相似三角形：角对应相等 $\Rightarrow$ 角亏相等 $\Rightarrow$ 面积相等。

(3) 面积相等而"相似比 $k \ne 1$"的三角形，边长差 $k$ 倍——但双曲几何有一条欧氏没有的刚性：\textbf{边长由角度决定}（把 (1) 反过来用：面积由角定，再由"双曲三角学"的边角关系，边也被角定死）。

(4) 于是相似比只能为 $1$：相似即全等。证毕。

对比一下你家窗户上贴的"缩小版三角形贴纸"——\textbf{欧氏世界允许缩放，双曲世界禁止缩放}。这道题的哲学分量：几何的公设不止决定"平行线有几条"，它渗进你能画什么图、能贴什么贴纸的每一个角落。
\end{shaonao}

\begin{chengshi}
"双曲几何自洽"的严格证明（在欧氏空间造模型——如庞加莱圆盘或贝尔特拉米伪球面局部实现）超出本书，但其意义必须说透：正因为有模型，双曲几何的\textbf{相对自洽性}成立——若它有矛盾，欧氏几何自己也会有矛盾。宇宙曲率 $\Omega_K = 0.001\pm0.002$ 是普朗克 2018 年的口径（结合 BAO 数据）；GPS 的 $38$ 微秒含狭义相对论（慢 $7\ \mu$s）与广义相对论（快 $45\ \mu$s）两种效应的叠加。黎曼演讲的原题《\"Uber die Hypothesen, welche der Geometrie zu Grunde liegen》是数学史上被引用最多的文本之一。
\end{chengshi}

篇五完。形状的世界我们已四管齐下：曲率（25）、测地线（26）、内在弯曲（27）、几何家族与宇宙（28）。最后一篇，把篇二的多项式重新请上台——这一次，\textbf{给每一个方程画一幅画}，一直画到密码学与费马大定理的大门口。
